%!TEX root = main.tex

\chapter{Limites e Continuidade}

\section{Sequências}

\begin{definition}
    \leavevmode
        \begin{enumerate}[leftmargin=*, align=left, label=\textbf{(\alph*)}]
            \item Uma sequência $(x_n)$ é \emph{convergente} e \emph{converge} para $a \in \R$ se para todo $\epsilon \in \R_{>0}$ existe $n_0 \in \N$ tal que, para todo $n \in \N$, se $n > n_0$, então $|x_n - a| < \epsilon$. Isso é denotado por
                \[
                    \lim x_n =a.
                \]
            Dizemos, nesse caso, que $a$ é o \emph{limite} de $(x_n)$.
            \item Uma sequência $(x_n)$ é \emph{divergente} se não é convergente.
        \end{enumerate}
\end{definition}

\begin{remark}
    As notações ``$\ds \lim_{n \to +\infty} x_n = a$'',  ``$\ds \lim_{n \in \N} x_n = a$'', ``$x_n \to a$ quando $n \to +\infty$'' e ``$x_n \to a$'' também são frequentemente usadas para indicar que $\lim x_n =a$.
\end{remark}

\begin{proposition}[Unicidade]
    O limite de uma sequência convergente é único.
\end{proposition}

\begin{proof}
    Provemos que se $(x_n)$ converge para $a \in \R$ e para $b \in \R$, então $a=b$. Suponha, por absurdo, que $a \neq b$, e seja $\epsilon := \dfrac{|b-a|}{2} \in \R_{>0}$. Como $\lim x_n = a$, existe $n_1 \in \N$ tal que, para todo $n \in \N$, se $n > n_1$, então $|x_n-a| < \epsilon$. Analogamente, como $\lim x_n = b$, existe $n_2 \in \N$ tal que, para todo $n \in \N$, se $n > n_2$, então $|x_n-b| < \epsilon$. Pois tome $n_0 := \max\{n_1,n_2\}$. Com isso, para todo $n \in \N$, se $n > n_0$, então $|x_n-a| < \epsilon$ e $|x_n-b| < \epsilon$, donde
        \begin{align*}
            |a-b| &= |(a-x_n) + (x_n-b)| \\
            &\leq |a-x_n| + |x_n-b| \\
            &< 2\epsilon = |a-b|,
        \end{align*}
    uma contradição. Logo, $a=b$. \blackproof
\end{proof}

\begin{theorem} \label{teo.calc1:algebra_das_sequencias}
    Sejam $(x_n)$ e $(y_n)$ convergentes.
        \begin{enumerate}[leftmargin=*, align=left, label=\textbf{(\alph*)}]
            \item A sequência $(x_n + y_n)$ é convergente e
                \[
                    \lim \, (x_n + y_n) = \lim x_n + \lim y_n.
                \]
            \item A sequência $(x_n \cdot y_n)$ é convergente e
                \[
                    \lim \, (x_n \cdot y_n) = \lim x_n \cdot \lim y_n.
                \]
            \item Se $\lim y_n \neq 0$, então a sequência $(x_n / y_n)$ é convergente e
                \[
                    \lim \left( \dfrac{x_n}{y_n} \right) = \dfrac{\lim x_n}{\lim y_n}.
                \]
        \end{enumerate}
\end{theorem}

\begin{proof}
    Sejam $\lim x_n = a$ e $\lim y_n = b$.
        \begin{enumerate}[leftmargin=*, align=left, label=\textbf{(\alph*)}]
            \item Seja $\epsilon \in \R_{>0}$. Como $\lim x_n = a$, para $\ds \epsilon_1 = \dfrac{\epsilon}{2} \in \R_{>0}$ existe $n_1 \in \N$ tal que, para todo $n \in \N$, se $n > n_1$, então $|x_n - a| < \epsilon_1$. Como $\lim y_n = b$, para $\ds \epsilon_2 = \dfrac{\epsilon}{2} \in \R_{>0}$ existe $n_2 \in \N$ tal que, para todo $n \in \N$, se $n > n_2$, então $|y_n - b| < \epsilon_2$. Pois tome $n_0 := \max \{n_1, n_2 \}$. Com isso, para todo $n \in \N$ tal que $n > n_0$, temos
                \begin{align*}
                    |(x_n + y_n) - (a + b)|
                        &\leq |x_n - a| + |y_n - b| \\
                        &< \dfrac{\epsilon}{2} + \dfrac{\epsilon}{2} \\
                        &= \epsilon,
                \end{align*}
            de modo que $\lim \, (x_n + y_n) = a + b$. \blackproof
            \item Seja $\epsilon \in \R_{>0}$. Como $\lim x_n = a$, para $\epsilon_1 = 1$ existe $n_1 \in \N$ tal que, para todo $n \in \N$, se $n > n_1$, então $|x_n - a| < \epsilon_1 = 1$; em particular, $|x_n| \leq 1 + |a|$. De fato, 
                \[
                    |x_n| = |x_n - a + a| \leq |x_n - a| + |a| < 1 + |a|.
                \]
            Como $\lim x_n = a$, para $\ds \epsilon_2 = \dfrac{\epsilon}{2(1 + |b|)} \in \R_{>0}$ existe $n_2 \in \N$ tal que, para todo $n \in \N$, se $n > n_2$, então $|x_n - a| < \epsilon_2$. Como $\lim y_n = b$, para $\epsilon_3 = \dfrac{\epsilon}{2(1 + |a|)} \in \R_{>0}$ existe $n_3 \in \N$ tal que, para todo $n \in \N$, se $n > n_3$, então $|y_n - b| < \epsilon_3$. Pois tome $n_0 := \max \{n_1, n_2, n_3 \}$. Observando que
                \[
                    x_n y_n - a b = x_n (y_n - b) + b (x_n - a),
                \]
            para todo $n \in \N$ tal que $n > n_0$, temos
                \begin{align*}
                    |x_n y_n - a b|
                        &\leq |x_n| |y_n - b| + |b| |x_n - a| \\
                        &< (1+|a|) \dfrac{\epsilon}{2(1+|a|)} + |b| \dfrac{\epsilon}{2(1+|b|)} \\
                        &< \dfrac{\epsilon}{2} + 1 \cdot \dfrac{\epsilon}{2} \\
                        &< \epsilon,
                \end{align*}
            de modo que $\lim \, (x_n \cdot y_n) = a \cdot b$. \blackproof
            \item Suponha $b \neq 0$. Seja $\epsilon \in \R_{>0}$. Como $\lim y_n = b$, para $\ds \epsilon_1 = \dfrac{|b|}{2} \in \R_{>0}$ existe $n_1 \in \N$ tal que, para todo $n \in \N$, se $n > n_1$, então $|y_n - b| < \epsilon_1$; em particular, $\dfrac{1}{|y_n|} < \dfrac{2}{|b|}$. De fato,
                \[
                    |y_n| = |y_n - b + b| \geq |b| - |y_n - b| > |b| - \dfrac{|b|}{2} = \dfrac{|b|}{2},
                \]
            donde $\dfrac{1}{|y_n|} < \dfrac{2}{|b|}$. Como $\lim x_n = a$, para $\ds \epsilon_2 = \dfrac{|b| \epsilon}{4} \in \R_{>0}$ existe $n_2 \in \N$ tal que, para todo $n \in \N$, se $n > n_2$, então $|x_n - a| < \epsilon_2$. Como $\lim y_n = b$, para $\ds \epsilon_3 = \dfrac{|b|^2 \epsilon}{4(1 + |a|)} \in \R_{>0}$ existe $n_3 \in \N$ tal que, para todo $n \in \N$, se $n > n_3$, então $|y_n - b| < \epsilon_3$. Pois tome $n_0 := \max \{n_1, n_2, n_3 \}$. Observando que
                \[
                    \dfrac{x_n}{y_n} - \dfrac{a}{b} = \dfrac{b(x_n - a) - a(y_n - b)}{y_n b},
                \]
            para todo $n \in \N$ tal que $n > n_0$, temos
                \begin{align*}
                    \left| \dfrac{x_n}{y_n} - \dfrac{a}{b} \right|
                        &= \dfrac{|b(x_n - a) - a(y_n - b)|}{|y_n| |b|} \\
                        &\leq \dfrac{|b| |x_n - a| + |a| |y_n - b|}{|y_n| |b|} \\
                        &< \dfrac{2}{|b|^2} \left( |b| \dfrac{|b| \epsilon}{4} + |a| \dfrac{|b|^2 \epsilon}{4(1 + |a|)} \right) \\
                        &< \dfrac{2}{|b|^2} \left( \dfrac{|b|^2 \epsilon}{4} + 1 \cdot \dfrac{|b|^2 \epsilon}{4} \right) \\
                        &= \epsilon,
                \end{align*}
            de modo que $\lim \left( \dfrac{x_n}{y_n} \right) = \dfrac{a}{b}$. \blackproof
        \end{enumerate}
\end{proof}

\begin{proposition}
    Seja $(x_n)$ convergente.
        \begin{enumerate}[leftmargin=*, align=left, label=\textbf{(\alph*)}]
            \item Para todo $a \in \R$, se $a < \lim x_n$, então, para todo $n$ suficientemente grande, $a < x_n$.
            \item Para todo $b \in \R$, se $\lim x_n < b$, então, para todo $n$ suficientemente grande, $x_n < b$.
        \end{enumerate}
\end{proposition}

\begin{proof}
    \leavevmode
        \begin{enumerate}[leftmargin=*, align=left, label=\textbf{(\alph*)}]
            \item Seja $\epsilon := \lim x_n - a \in \R_{>0}$. Então existe $n_0 \in \N$ tal que, para todo $n \in \N$, se $n > n_0$, então $\lim x_n - \epsilon < x_n < \lim x_n + \epsilon$; como $a = \lim x_n - \epsilon$, temos $a < x_n$ para todo $n$ suficientemente grande. \blackproof
            \item Seja $\epsilon := b - \lim x_n \in \R_{>0}$. Então existe $n_0 \in \N$ tal que, para todo $n \in \N$, se $n > n_0$, então $\lim x_n - \epsilon < x_n < \lim x_n + \epsilon$; como $b = \lim x_n + \epsilon$, temos $x_n < b$ para todo $n$ suficientemente grande. \blackproof
        \end{enumerate}
\end{proof}

\begin{corollary}
    Sejam $(x_n)$ e $(y_n)$ convergentes. Se $x_n \leq y_n$ para todo $n$ suficientemente grande, então $\lim x_n \leq \lim y_n$.
\end{corollary}

\begin{proof}
    Sejam $a := \lim x_n$ e $b := \lim y_n$. Suponha, por absurdo, que $b < a$. Seja $c \in \R$ tal que $b < c < a$. Então, para todo $n$ suficientemente grande, teríamos $y_n < c < x_n$, o que contradiz a hipótese. \blackproof
\end{proof}

\begin{theorem}[Confronto]
    Sejam $(x_n)$, $(y_n)$ e $(z_n)$ tais que $x_n \leq y_n \leq z_n$ para todo $n$ suficientemente grande. Se $\lim x_n = a$ e $\lim z_n = a$, então $\lim y_n = a$.
\end{theorem}

\begin{proof}
    Seja $\epsilon \in \R_{>0}$. Como $\lim x_n = a$, existe $n_1 \in \N$ tal que, para todo $n \in \N$, se $n > n_1$, então $a-\epsilon < x_n < a+\epsilon$. Como $\lim z_n = a$, existe $n_2 \in \N$ tal que, para todo $n \in \N$, se $n > n_2$, então $a-\epsilon < z_n < a+\epsilon$. Como $x_n \leq y_n \leq z_n$ para todo $n$ suficientemente grande, existe $n_3 \in \N$ tal que, para todo $n \in \N$, se $n > n_3$, então $x_n \leq y_n \leq z_n$. Seja $n_0 := \max \{n_1, n_2, n_3 \}$. Com isso, para todo $n \in \N$, se $n > n_0$, então 
        \[
            a-\epsilon < x_n \leq y_n \leq z_n < a+\epsilon,
        \]
    de modo que $|y_n - a| < \epsilon$. Logo, $\lim y_n = a$. \blackproof
\end{proof}

\begin{proposition}
    Toda sequência convergente é limitada.
\end{proposition}

\begin{proof}
    Sejam $(x_n)$ e $a \in \R$ tais que $\lim x_n = a$. Para $\epsilon = 1$, existe $n_0 \in \N$ tal que, para todo $n \in \N$, se $n > n_0$, então $|x_n - a| < 1$, isto é, $a-1 < x_n < a+1$. Então, sendo $S := \{x_1, x_2, \ldots, x_{n_0}, a-1, a+1\}$, temos $|x_n| \leq \max\{|\max\{S\}|, |\min\{S\}|\}$ para todo $n \in \N$, de modo que $(x_n)$ é limitada. \blackproof
\end{proof}

\begin{theorem}[Convergência Monótona] \label{teo.calc1:convergencia_monotona}
    \leavevmode
        \begin{enumerate}[leftmargin=*, align=left, label=\textbf{(\alph*)}]
            \item Toda sequência crescente e limitada superiormente é convergente.
            \item Toda sequência decrescente e limitada inferiormente é convergente.
            \item Toda sequência monótona e limitada é convergente.
        \end{enumerate}
\end{theorem}

\begin{proof}
    \leavevmode
        \begin{enumerate}[leftmargin=*, align=left, label=\textbf{(\alph*)}]
            \item Seja $(x_n)$ crescente e limitada superiormente. Como $\Im{x}$ é não vazio e limitado superiormente, existe $\sup{\Im{x}}$. Agora, seja $\epsilon \in \R_{>0}$. Então existe $n_0 \in \N$ tal que $\sup{\Im{x}} - \epsilon < x_{n_0}$ pois $\sup{\Im{x}} - \epsilon$ não é uma cota superior de $\Im{x}$. Como $(x_n)$ é crescente, para todo $n \in \N$, se $n > n_0$, então
                \[
                    \sup{\Im{x}} - \epsilon < x_{n_0} \leq x_n \leq \sup{\Im{x}} < \sup{\Im{x}} + \epsilon,
                \]
            de modo que $\lim x_n = \sup{\Im{x}}$. \blackproof
            \item Seja $(x_n)$ decrescente e limitada inferiormente. Como $\Im{x}$ é não vazio e limitado inferiormente, existe $\inf{\Im{x}}$. Agora, seja $\epsilon \in \R_{>0}$. Então existe $n_0 \in \N$ tal que $x_{n_0} < \inf{\Im{x}} + \epsilon$ pois $\inf{\Im{x}} + \epsilon$ não é uma cota inferior de $\Im{x}$. Como $(x_n)$ é decrescente, para todo $n \in \N$, se $n > n_0$, então
                \[
                    \inf{\Im{x}} - \epsilon < \inf{\Im{x}} \leq x_n \leq x_{n_0} < \inf{\Im{x}} + \epsilon,
                \]
            de modo que $\lim x_n = \inf{\Im{x}}$. \blackproof
            \item Segue dos itens acima. \blackproof
        \end{enumerate}
\end{proof}

\subsection{Subsequências}

\begin{definition} \label{defi:subseq}
    Sejam $x : \N \to \R$ e $n : \N \to \N$ estritamente crescente. Dizemos que $x \circ n : \N \to \R$ é uma \emph{subsequência} de $x$. Isso é denotado por $(x_{n_{k}})$.
\end{definition}

\begin{proposition} \fonte[26]{analisereal1}
    Seja $(x_n)$ monótona. Se existe $(x_{n_{k}})$ limitada, então $(x_n)$ é limitada.
\end{proposition}

\begin{proof}
    \blackproof
\end{proof}

\begin{proposition} \label{prop.calc1:subseqconv}
    Se $(x_n)$ converge para $a \in \R$, então toda subsequência $(x_{n_{k}})$ converge para $a$.
\end{proposition}

\begin{proof}
    Seja $\epsilon \in \R_{>0}$. Como $\lim x_n = a$, existe $n_0 \in \N$ tal que, para todo $n \in \N$, se $n > n_0$, então $|x_n - a| < \epsilon$. Como $(n_k)$ é estritamente crescente, existe $k_0 \in \N$ tal que $n_{k_0} > n_0$; com isso, para todo $k \in \N$, se $k > k_0$, então $n_k > n_{k_0} > n_0$, de modo que $|x_{n_{k}} - a| < \epsilon$. Logo, $\lim x_{n_{k}} = a$. \blackproof
\end{proof}

\begin{theorem}[Bolzano-Weierstrass] \label{teo.calc1:bolzanoweierstrass}
    Toda sequência limitada possui uma subsequência convergente.
\end{theorem}

\begin{proof}
    \blackproof
\end{proof}

\section{Séries}

\section{Topologia da Reta}

\begin{definition}
    Sejam $a \in \R$ e $r \in \R_{>0}$.
        \begin{enumerate}[leftmargin=*, align=left, label=\textbf{(\alph*)}]
            \item A \emph{$r$-vizinhança de $a$} é definida como
                \[
                    V_r(a) := \{x \in \R : |x-a|<r \}.
                \]
            \item A \emph{$r$-vizinhança furada de $a$} é definida como
                \[
                    V^*_r(a) := \{x \in \R : 0<|x-a|<r \}.
                \]
        \end{enumerate}
\end{definition}

\begin{corollary}
    Sejam $a \in \R$ e $r \in \R_{>0}$.
        \begin{enumerate}[leftmargin=*, align=left, label=\textbf{(\alph*)}]
            \item $V_r(a) = \left]a-r, a+r \right[$.
            \item $V^*_r(a) = V_r(a) \setminus \{a\}$.
        \end{enumerate}
\end{corollary}

\begin{proof}
    \leavevmode
        \begin{enumerate}[leftmargin=*, align=left, label=\textbf{(\alph*)}]
            \item \blackproof
            \item \blackproof
        \end{enumerate}
\end{proof}

\begin{definition}
    Dizemos que $a \in \R$ é um \emph{ponto de acumulação} de $A \subseteq \R$ se 
        \[
            V^*_{r}(a) \cap A \neq \emptyset
        \]
    para todo $r \in \R_{>0}$. O conjunto de todos os pontos de acumulação de $A$ é denotado por $A'$.
\end{definition}

\begin{theorem} \label{teo.calc1:a_linha_sse_sequencia}
    Sejam $A \subseteq \R$ e $a \in \R$. Então, $a \in A'$ se, e somente se, existe $(x_n)$ em $A_{\neq a}$ tal que $\lim x_n = a$. 
\end{theorem}

\begin{proof}
    Ver \cite{analisereal1}, página 49, teorema 6. \blackproof
\end{proof}

\section{Limites}

\begin{definition}
    Dizemos que $L \in \R$ é o \emph{limite} de uma função $f : A \subseteq \R \to \R$ em $a \in A'$ se para todo $\epsilon \in \R_{>0}$ existe $\delta = \delta(\epsilon, a) \in \R_{>0}$ tal que, para todo $x \in A$, se $0<|x-a|<\delta$, então $|f(x)-L| < \epsilon$. Isso é denotado por 
        \[
            \lim_{x \to a} f(x) = L.
        \]
\end{definition}

\begin{lemma} \label{calc1.lemma:existe_x0}
    Sejam $a \in \R$ e $A \subseteq \R$. Temos $a \in A'$ se, e somente se, para todo $\delta \in \R_{>0}$ existe $x_0 \in V^*_{\delta}(a) \cap A$.
\end{lemma}

\begin{proof}
    Por um lado ($\Rightarrow$), suponha que $a \in A'$. Seja $\delta \in \R_{>0}$. Como $a \in A'$, temos $V^*_{\delta}(a) \cap A \neq \emptyset$. Logo, existe $x_0 \in V^*_{\delta}(a) \cap A$. Por outro lado ($\Leftarrow$), suponha que, para todo $\delta \in \R_{>0}$, existe $x_0 \in V^*_{\delta}(a) \cap A$. Então $V^*_{\delta}(a) \cap A \neq \emptyset$ para todo $\delta \in \R_{>0}$, de modo que $a \in A'$. \blackproof
\end{proof}

\begin{proposition}[Unicidade] \label{teo.calc1:unicidade_limite_bilateral}
    O limite de $f : A \subseteq \R \to \R$ em $a \in A'$, quando existe, é único.
\end{proposition}

\begin{proof} \fonte[135]{cursodeanalise1}
    Provemos que se $\ds \lim_{x \to a} f(x) = L_1$ e $\ds \lim_{x \to a} f(x) = L_2$, então $L_1 = L_2$. Suponha, por absurdo, que $L_1 \neq L_2$, e seja $\epsilon := \dfrac{|L_1 - L_2|}{2} \in \R_{>0}$. Como $\ds \lim_{x \to a} f(x) = L_1$, existe $\delta_1 \in \R_{>0}$ tal que, para todo $x \in A$, se $0 < |x-a| < \delta_1$, então $|f(x) - L_1| < \epsilon$. Analogamente, como $\ds \lim_{x \to a} f(x) = L_2$, existe $\delta_2 \in \R_{>0}$ tal que, para todo $x \in A$, se $0 < |x-a| < \delta_2$, então $|f(x) - L_2| < \epsilon$. Pois tome $\delta := \min\{\delta_1, \delta_2 \}$. como $a \in A'$, pelo \cref{calc1.lemma:existe_x0} existe $x_0 \in A$ tal que $0<|x_0-a|<\delta$; com isso, obtemos $|f(x_0) - L_1| < \epsilon$ e $|f(x_0) - L_2| < \epsilon$, donde
        \begin{align*}
            |L_1 - L_2| &= |(f(x_0) - L_2) - (f(x_0)-L_1)| \\
            &\leq |f(x_0) - L_2| + |f(x_0)-L_1| \\
            &< 2 \epsilon = |L_1 - L_2|,
        \end{align*}
    uma contradição. Logo, $L_1 = L_2$. \blackproof
\end{proof}

\begin{theorem} \label{teo.calc1:limite_sequencial} \fonte[136]{cursodeanalise1}
    Sejam $f : A \subseteq \R \to \R$, $a \in A'$ e $L \in \R$. Temos $\ds \lim_{x \to a} f(x) = L$ se, e somente se, para toda sequência $(x_n)$ em $A \setminus \{a\}$, se $\lim x_n = a$, então $\lim f(x_n) = L$.
\end{theorem}

\begin{proof}
    \leavevmode
        \begin{itemize}[align=left]
            \item[($\Rightarrow$)] Suponha que $\ds \lim_{x \to a} f(x) = L$, e seja $(x_n)$ em $A \setminus \{a\}$ tal que $\lim x_n = a$. Seja $\epsilon \in \R_{>0}$. Então existe $\delta \in \R_{>0}$ tal que, para todo $x \in A$, se $0<|x-a|<\delta$, então $|f(x) - L|<\epsilon$. Como $\lim x_n = a$, existe $N \in \N$ tal que, para todo $n \in \N$, se $n > N$, então $|x_n - a| < \delta$. Como $x_n \neq a$ para todo $n \in \N$, temos $0<|x_n - a| < \delta$ para todo $n > N$. Logo, para todo $n \in \N$, se $n > N$, então $|f(x_n) - L| < \epsilon$, isto é, $\lim f(x_n) = L$.
            \item[($\Leftarrow$)] Suponha que $\lim f(x_n) = L$ para toda sequência $(x_n)$ em $A \setminus \{a\}$ tal que $\lim x_n = a$. Suponha, por absurdo, que $\ds \lim_{x \to a} f(x) \neq L$. Então existe $\xi \in \R_{>0}$ tal que, para todo $\delta \in \R_{>0}$, existe $x = x(\delta) \in A$ tal que $0<|x(\delta)-a|<\delta$ e $|f(x(\delta))-L| \geq \xi$. Em particular, para todo $n \in \N$, tomando $\delta_n := 1/n$, existe $x_n := x(\delta_n) \in A$ tal que $0<|x_n - a|< 1/n$ e $|f(x_n)-L| \geq \xi$.
                \begin{claim}
                    A sequência $(x_n)$ converge para $a$, isto é, $\lim x_n = a$.
                \end{claim}
                \begin{midproof}
                    Seja $\epsilon \in \R_{>0}$. Pela propriedade arquimediana, existe $N \in \N$ tal que $1/N < \epsilon$. Então, para todo $n \in \N$, se $n > N$, temos $1/n < 1/N$, de modo que
                        \[
                            |x_n - a| < 1/n < 1/N < \epsilon,
                        \]
                    isto é, $\lim x_n = a$ e $(x_n)$ converge para $a$. \whiteproof
                \end{midproof}
            Como $0<|x_n - a|$ para todo $n \in \N$, então $(x_n)$ é uma sequência em $A \setminus \{a\}$. Como $\lim x_n = a$, deveríamos ter $\lim f(x_n) = L$, mas existe $\xi \in \R_{>0}$ tal que $|f(x_n)-L| \geq \xi$ para todo $n \in \N$, uma contradição. Logo, $\ds \lim_{x \to a} f(x) = L$.
        \end{itemize}
    Isso completa a prova. \blackproof
\end{proof}

\begin{theorem} \label{teo.calc1:algebra_dos_limites}
    Sejam $f,g : A \subseteq \R \to \R$ e $a \in A'$ tais que existem $\ds \lim_{x \to a} f(x)$ e $\ds \lim_{x \to a} g(x)$.
        \begin{enumerate}[leftmargin=*, align=left, label=\textbf{(\alph*)}]
            \item O limite da soma é a soma dos limites:
                \[
                    \lim_{x \to a} \, (f+g)(x) = \lim_{x \to a} f(x) + \lim_{x \to a} g(x).
                \]
            \item O limite do produto é o produto dos limites:
                \[
                    \lim_{x \to a} \, (f \cdot g)(x) = \lim_{x \to a} f(x) \cdot \lim_{x \to a} g(x).
                \]
            \item O limite do quociente é igual ao quociente dos limites, desde que o denominador seja diferente de $0$: se $\ds \lim_{x \to a} g(x) \neq 0$, então
                \[
                    \lim_{x \to a} \left(\dfrac{f}{g} \right)(x) = \dfrac{\ds \lim_{x \to a} f(x)}{\ds \lim_{x \to a} g(x)}.
                \]
        \end{enumerate}
\end{theorem}

\begin{proof}
    Sejam $\ds \lim_{x \to a} f(x) = L_1$ e $\ds \lim_{x \to a} g(x) = L_2$.
        \begin{enumerate}[leftmargin=*, align=left, label=\textbf{(\alph*)}]
            \item Seja $\epsilon \in \R_{>0}$. Como $\ds \lim_{x \to a} f(x) = L_1$, para $\ds \epsilon_1 = \dfrac{\epsilon}{2} \in \mathbb{R}_{>0}$ existe $\delta_1 \in \mathbb{R}_{>0}$ tal que, para todo $x \in A$, se $0<|x-a|<\delta_1$, então $|f(x) - L_1| < \epsilon_1$. Como $\ds \lim_{x \to a} g(x) = L_2$, para $\ds \epsilon_2 = \dfrac{\epsilon}{2} \in \mathbb{R}_{>0}$ existe $\delta_2 \in \mathbb{R}_{>0}$ tal que, para todo $x \in A$, se $0<|x-a|<\delta_2$, então $|g(x) - L_2| < \epsilon_2$. Pois tome $\delta := \min \{\delta_1, \delta_2 \}$. Com isso, para todo $x \in A$ tal que $0 < |x-a| < \delta$, temos
                \begin{align*}
                    |(f(x) + g(x)) - (L_1 + L_2)|
                        &\leq |f(x) - L_1| + |g(x) - L_2| \\
                        &< \dfrac{\epsilon}{2} + \dfrac{\epsilon}{2} \\
                        &= \epsilon,
                \end{align*}
            de modo que $\ds \lim_{x \to a} (f+g)(x) = L_1 + L_2$. \blackproof
            \item \fonte[30]{aurichi}
            Seja $\epsilon \in \R_{>0}$. Como $\ds \lim_{x \to a} f(x) = L_1$, para $\epsilon_1 = 1$ existe $\delta_1 \in \R_{>0}$ tal que, para todo $x \in A$, se $0<|x-a|<\delta_1$, então $|f(x) - L_1| < \epsilon_1 = 1$; em particular, $|f(x)| \leq 1 + |L_1|$. De fato, 
                \[
                    |f(x)| = |f(x) - L_1 + L_1| \leq |f(x) - L_1| + |L_1| < 1 + |L_1|.
                \]
            Como $\ds \lim_{x \to a} f(x) = L_1$, para $\ds \epsilon_2 = \dfrac{\epsilon}{2(1 + |L_2|)} \in \R_{>0}$ existe $\delta_2 \in \R_{>0}$ tal que, para todo $x \in A$, se $0<|x-a|<\delta_2$, então $|f(x) - L_1| < \epsilon_2$. Como $\ds \lim_{x \to a} g(x) = L_2$, para $\epsilon_3 = \dfrac{\epsilon}{2(1 + |L_1|)} \in \R_{>0}$ existe $\delta_3 \in \R_{>0}$ tal que, para todo $x \in A$, se $0<|x-a|<\delta_3$, então $|g(x) - L_2| < \epsilon_3$. Pois tome $\delta := \min \{\delta_1, \delta_2, \delta_3 \}$. Observando que
                \[
                    f(x) g(x) - L_1 L_2 = f(x) (g(x) - L_2) + L_2 (f(x) - L_1),
                \]
            para todo $x \in A$ tal que $0 < |x-a| < \delta$, temos
                \begin{align*}
                    |f(x) g(x) - L_1 L_2|
                        &\leq |f(x)| |g(x) - L_2| + |L_2| |f(x) - L_1| \\
                        &< (1+|L_1|) \dfrac{\epsilon}{2(1+|L_1|)} + |L_2| \dfrac{\epsilon}{2(1+|L_2|)} \\
                        &< \dfrac{\epsilon}{2} + 1 \cdot \dfrac{\epsilon}{2} \\
                        &= \epsilon,
                \end{align*}
            de modo que $\ds \lim_{x \to a} f(x)g(x) = L_1 L_2$. \blackproof
            \item Suponha $L_2 \neq 0$. Seja $\epsilon \in \R_{>0}$. Como $\ds \lim_{x \to a} g(x) = L_2$, para $\ds \epsilon_1 = \dfrac{|L_2|}{2} \in \R_{>0}$ existe $\delta_1 \in \R_{>0}$ tal que, para todo $x \in A$, se $0<|x-a|<\delta_1$, então $|g(x) - L_2| < \epsilon_1$; em particular, $\dfrac{1}{|g(x)|} < \dfrac{2}{|L_2|}$. De fato,
                \[
                    |g(x)| = |g(x) - L_2 + L_2| \geq |L_2| - |g(x) - L_2| > |L_2| - \dfrac{|L_2|}{2} = \dfrac{|L_2|}{2},
                \]
            donde $\dfrac{1}{|g(x)|} < \dfrac{2}{|L_2|}$. Como $\ds \lim_{x \to a} f(x) = L_1$, para $\ds \epsilon_2 = \dfrac{|L_2| \epsilon}{4} \in \R_{>0}$ existe $\delta_2 \in \R_{>0}$ tal que, para todo $x \in A$, se $0<|x-a|<\delta_2$, então $|f(x) - L_1| < \epsilon_2$. Como $\ds \lim_{x \to a} g(x) = L_2$, para $\ds \epsilon_3 = \dfrac{|L_2|^2 \epsilon}{4(1 + |L_1|)} \in \R_{>0}$ existe $\delta_3 \in \R_{>0}$ tal que, para todo $x \in A$, se $0<|x-a|<\delta_3$, então $|g(x) - L_2| < \epsilon_3$. Pois tome $\delta := \min \{\delta_1, \delta_2, \delta_3 \}$. Observando que
                \[
                    \dfrac{f(x)}{g(x)} - \dfrac{L_1}{L_2} = \dfrac{L_2(f(x) - L_1) - L_1(g(x) - L_2)}{g(x)L_2},
                \]
            para todo $x \in A$ tal que $0 < |x-a| < \delta$, temos
                \begin{align*}
                    \left| \dfrac{f(x)}{g(x)} - \dfrac{L_1}{L_2} \right|
                        &= \dfrac{|L_2(f(x) - L_1) - L_1(g(x) - L_2)|}{|g(x)| |L_2|} \\
                        &\leq \dfrac{|L_2| |f(x) - L_1| + |L_1| |g(x) - L_2|}{|g(x)| |L_2|} \\
                        &< \dfrac{2}{|L_2|^2} \left( |L_2| \dfrac{|L_2| \epsilon}{4} + |L_1| \dfrac{|L_2|^2 \epsilon}{4(1 + |L_1|)} \right) \\
                        &< \dfrac{2}{|L_2|^2} \left( \dfrac{|L_2|^2 \epsilon}{4} + 1 \cdot \dfrac{|L_2|^2 \epsilon}{4} \right) \\
                        &= \epsilon,
                \end{align*}
            de modo que $\ds \lim_{x \to a} \dfrac{f(x)}{g(x)} = \dfrac{L_1}{L_2}$. \blackproof
        \end{enumerate}
\end{proof}

\begin{corollary}
    Sejam $n \in \N$, $A \subseteq \R$, $a \in A'$. Para cada $i \in [n]$, seja $f_i : A \to \R$ tal que existe $\ds \lim_{x \to a} f_i(x)$.
        \begin{enumerate}[leftmargin=*, align=left, label=\textbf{(\alph*)}]
            \item O limite da soma é igual à soma dos limites:
                \[
                    \lim_{x \to a} \left( \sum_{i=1}^{n} f_i(x) \right) = \sum_{i=1}^{n} \left( \lim_{x \to a} f_i(x) \right).
                \]
            \item O limite do produto é igual ao produto dos limites:
                \[
                    \lim_{x \to a} \left( \prod_{i=1}^{n} f_i(x) \right) = \prod_{i=1}^{n} \left( \lim_{x \to a} f_i(x) \right).
                \]
        \end{enumerate}
\end{corollary}

\begin{proof}
    Por indução. \blackproof
        \begin{enumerate}[leftmargin=*, align=left, label=\textbf{(\alph*)}]
            \item \blackproof
            \item \blackproof
        \end{enumerate}
\end{proof}

\begin{theorem} \fonte[135]{cursodeanalise1}
    Sejam $f : A \subseteq \R \to \R$, $B \subseteq A$, $a \in B'$ e $L \in \R$.
        \begin{enumerate}[leftmargin=*, align=left, label=\textbf{(\alph*)}]
            \item Se $\ds \lim_{x \to a} f(x) = L$, então $\ds \lim_{x \to a} f \restriction_B(x) = L$.
            \item Se $B = I \cap A$, onde $I \subseteq \R$ é um intervalo aberto tal que $a \in I$, então
                \[
                    \lim_{x \to a} f \restriction_{B}(x) = L \Rightarrow \lim_{x \to a} f(x) = L
                \]
        \end{enumerate}
\end{theorem}

\begin{proof}
    Seja $\epsilon \in \R_{>0}$.
        \begin{enumerate}[leftmargin=*, align=left, label=\textbf{(\alph*)}]
            \item Como $\ds \lim_{x \to a} f(x) = L$, existe $\delta \in \R_{>0}$ tal que, para todo $x \in A$, se $0<|x-a|<\delta$, então $|f(x)-L|<\epsilon$. Como $B \subseteq A$, para todo $x \in B$, se $0<|x-a|<\delta$, então $x \in A$ e, portanto,
                \[
                    |f \restriction_B(x)-L| = |f(x)-L| < \epsilon.
                \]
            Logo, $\ds \lim_{x \to a} f \restriction_B(x) = L$. \blackproof
            \item Como $\ds \lim_{x \to a} f \restriction_B(x) = L$, existe $\delta_1 \in \R_{>0}$ tal que, para todo $x \in B$, se $0<|x-a|<\delta_1$, então $|f \restriction_B(x)-L|<\epsilon$. Como $I$ é aberto e $a \in I$, existe $\delta_2 \in \R_{>0}$ tal que $V_{\delta_2}(a) \subseteq I$. Pois tome $\delta := \min \{\delta_1, \delta_2\}$. Logo, para todo $x \in A$, se $0<|x-a|<\delta$, então $x \in V_{\delta_2}(a) \subseteq I$, donde $x \in I \cap A = B$; com isso,
                \[
                    |f(x)-L| = |f \restriction_B(x)-L| < \epsilon.
                \]
            Logo, $\ds \lim_{x \to a} f(x) = L$. \blackproof
        \end{enumerate}
\end{proof}

\begin{theorem}[Conservação da Ordem] \label{teo.calc1:conservacao_da_ordem}
    Sejam $f,g : A \subseteq \R \to \R$ e $a \in A'$.
        \begin{enumerate}[leftmargin=*, align=left, label=\textbf{(\alph*)}]
            \item Se $\ds \lim_{x \to a} f(x) < \lim_{x \to a} g(x)$, então existe $r \in \R_{>0}$ tal que $f < g$ em $V^*_r(a) \cap A$.
            \item Se existe $r \in \R_{>0}$ tal que $f \leq g$ em $V^*_r(a) \cap A$, então $\ds \lim_{x \to a} f(x) \leq \lim_{x \to a} g(x)$.
        \end{enumerate}
\end{theorem}

\begin{proof}
    Sejam $\ds \lim_{x \to a} f(x) = L_1$ e $\ds \lim_{x \to a} g(x) = L_2$.
        \begin{enumerate}[leftmargin=*, align=left, label=\textbf{(\alph*)}]
            \item Seja $\ds \epsilon = \dfrac{L_2 - L_1}{2} \in \R_{>0}$. Como $\ds \lim_{x \to a} f(x) = L_1$, existe $\delta_1 \in \R_{>0}$ tal que, para todo $x \in A$, se $0<|x-a|<\delta_1$, então $|f(x) - L_1| < \epsilon$; em particular, $f(x) < L_1 + \epsilon$. Como $\ds \lim_{x \to a} g(x) = L_2$, existe $\delta_2 \in \R_{>0}$ tal que, para todo $x \in A$, se $0<|x-a|<\delta_2$, então $|g(x) - L_2| < \epsilon$; em particular, $L_2 - \epsilon < g(x)$. Pois tome $r := \min \{\delta_1, \delta_2 \}$. Logo, para todo $x \in V^*_r(a) \cap A$, isto é, $x \in A$ tal que $0 < |x-a| < r$, temos
                \[
                    f(x) < L_1 + \epsilon = L_2 - \epsilon < g(x),
                \]
            isto é, $f(x) < g(x)$ para todo $x \in V^*_r(a) \cap A$. \blackproof
            \item Suponha, por absurdo, que $L_2 < L_1$. Pelo item anterior, existe $r' \in \R_{>0}$ tal que $f > g$ em $V^*_{r'}(a) \cap A$. Pois tome $\delta := \min \{r, r'\}$. Então $f \leq g$ e $f > g$ em $V^*_{\delta}(a) \cap A$, uma contradição. Logo, $L_1 \leq L_2$. \blackproof
            \repeatitem Suponha, por absurdo, que $L_2 < L_1$. Como $\ds \lim_{x \to a} (f(x) - g(x)) = L_1 - L_2$, tomando $L:= L_1 - L_2$, para $\epsilon = L/2 \in \R_{>0}$ existe $\delta_1 \in \R_{>0}$ tal que, para todo $x \in A$, se $0<|x-a|<\delta_1$, então $|(f(x) - g(x)) - L| < L/2$. Em particular, $g < f$ em $V^*_{\delta_1}(a) \cap A$. Agora, seja $\delta := \min \{r, \delta_1\}$. Como $a \in A'$, existe $x_0 \in V^*_{\delta}(a) \cap A$. Como $x_0 \in V^*_{\delta_1}(a) \cap A$, temos $g(x_0) < f(x_0)$, e como $x_0 \in V^*_{r}(a) \cap A$, temos $f(x_0) \leq g(x_0)$, uma contradição. Logo, $L_1 \leq L_2$. \blackproof
        \end{enumerate}
\end{proof}

\begin{comment}
\begin{theorem}[Comparação]
    Sejam $f,g : A \subseteq \R \to \R$ e $a \in A'$. Se existe $r \in \R_{>0}$ tal que $f \leq g$ em $V^*_r(a) \cap A$, então
        \[
            \lim_{x \to a} f(x) \leq \lim_{x \to a} g(x),
        \]
    desde que os limites existam.
    %tais que $f(x) \leq g(x)$ para todo $x \in A$. Se existem $\ds \lim_{x \to a} f(x)$ e $\ds \lim_{x \to a} g(x)$, então $\ds \lim_{x \to a} f(x) \leq \lim_{x \to a} g(x)$.
    %Se $\ds \lim_{x \to a} f(x) < \lim_{x \to a} g(x)$, então existe $r \in \R_{>0}$ tal que $f \leq g$ em $V^*_{r}(a) \cap A$.
\end{theorem}

\begin{proof}
    Sejam $\ds \lim_{x \to a} f(x) = L_1$ e $\ds \lim_{x \to a} g(x) = L_2$. Suponha, por absurdo, que $L_2 < L_1$. Como
        $ \ds
            \lim_{x \to a} (f(x) - g(x)) = L_1 - L_2,
        $
    tomando $L:= L_1 - L_2$, para $\epsilon = \frac{L}{2} \in \R_{>0}$ existe $\delta_1 \in \R_{>0}$ tal que, para todo $x \in A$, se $0<|x-a|<\delta_1$, então
        $
            |(f(x) - g(x)) - L| < \frac{L}{2}.
        $
    Em particular,
    %$\ds 0 < \dfrac{L_1 - L_2}{2} < (f(x) - g(x)),$ de modo que
    $g(x) < f(x)$ para todo $x \in V^*_{\delta_1}(a) \cap A$. Agora, seja $\delta := \min \{r, \delta_1\}$. Como $a \in A'$, existe $x_0 \in V^*_{\delta}(a) \cap A$. Como $x_0 \in V^*_{\delta_1}(a) \cap A$, temos $g(x_0) < f(x_0)$, e como $x_0 \in V^*_{r}(a) \cap A$, temos $f(x_0) \leq g(x_0)$, uma contradição. Logo $L_1 \leq L_2$. \blackproof
\end{proof}

\begin{proof}
    Sejam $\ds \lim_{x \to a} f(x) = L_1$ e $\ds \lim_{x \to a} g(x) = L_2$. Suponha, por absurdo, que $L_2 < L_1$. Como
        \[
            \lim_{x \to a} (f(x) - g(x)) = \lim_{x \to a} f(x) - \lim_{x \to a} g(x) = L_1 - L_2,
        \]
    para $\epsilon = \dfrac{L_1 - L_2}{2} \in \R_{>0}$ existe $\delta_1 \in \R_{>0}$ tal que, para todo $x \in A$, se $0<|x-a|<\delta_1$, então
        \[
            |(f(x) - g(x)) - (L_1 - L_2)| < \dfrac{L_1 - L_2}{2},
        \]
    isto é,
        \[
            -\dfrac{L_1 - L_2}{2} < (f(x) - g(x)) - (L_1 - L_2) < \dfrac{L_1 - L_2}{2}.
        \]
    Em particular,
    %$\ds 0 < \dfrac{L_1 - L_2}{2} < (f(x) - g(x)),$ de modo que
    $g(x) < f(x)$ para todo $x \in V^*_{\delta_1}(a) \cap A$. Agora, seja $\delta := \min \{r, \delta_1\}$. Como $a \in A'$, existe $x_0 \in V^*_{\delta}(a) \cap A$. Como $x_0 \in V^*_{\delta_1}(a) \cap A$, temos $g(x_0) < f(x_0)$, e como $x_0 \in V^*_{r}(a) \cap A$, temos $f(x_0) \leq g(x_0)$, uma contradição. Logo $L_1 \leq L_2$. \blackproof
\end{proof}
\end{comment}

\begin{theorem}[Confronto]
    Sejam $f,g,h : A \subseteq \R \to \R$, $a \in A'$ e $L \in \R$. Se existe $r \in \R_{>0}$ tal que $f \leq g \leq h$ em $V^*_r(a) \cap A$, e se $\ds \lim_{x \to a} f(x) = L$ e $\ds \lim_{x \to a} h(x) = L$, então $\ds \lim_{x \to a} g(x) = L$.
\end{theorem}

\begin{proof}
    Seja $\epsilon \in \R_{>0}$. Como $\ds \lim_{x \to a} f(x) = L$, existe $\delta_1 \in \R_{>0}$ tal que, para todo $x \in A$, se $0<|x-a|<\delta_1$, então $L-\epsilon < f(x) < L+\epsilon$. Como $\ds \lim_{x \to a} h(x) = L$, existe $\delta_2 \in \R_{>0}$ tal que, para todo $x \in A$, se $0<|x-a|<\delta_2$, então $L-\epsilon < h(x) < L+\epsilon$. Seja $\delta := \min \{r, \delta_1, \delta_2 \}$. Com isso, para todo $x \in A$, se $0<|x-a|<\delta$, então 
        \[
            L-\epsilon < f(x) \leq g(x) \leq h(x) < L+\epsilon,
        \]
    isto é, $|g(x) - L| < \epsilon$. Logo, $\ds \lim_{x \to a} g(x) = L$. \blackproof
\end{proof}

\begin{definition}
    Dizemos que $f : A \subseteq \R \to \R$ é \emph{localmente limitada} em $a \in A \cup A'$ se existe $r \in \R_{>0}$ tal que $f \restriction_{V_r(a) \cap A}$ é limitada.
\end{definition}

\begin{proposition}
    Sejam $f,g : A \subseteq \R \to \R$ e $a \in A'$.
        \begin{enumerate}[leftmargin=*, align=left, label=\textbf{(\alph*)}]
            \item Se $\ds \lim_{x \to a} f(x) \in \R$, então $f$ é localmente limitada em $a$.
            \item Se $\ds \lim_{x \to a} f(x) = 0$ e $g$ é localmente limitada em $a$, então $\ds \lim_{x \to a} f(x)g(x) = 0$.
        \end{enumerate}
\end{proposition}

\begin{proof}
    \leavevmode
        \begin{enumerate}[leftmargin=*, align=left, label=\textbf{(\alph*)}]
            \item Seja $\ds L := \lim_{x \to a} f(x)$. Para $\epsilon = 1$, existe $\delta \in \R_{>0}$ tal que, para todo $x \in A$, se $0 < |x-a| < \delta$, então $|f(x) - L| < 1$. Com isso, para todo $x \in V^*_{\delta}(a) \cap A$, temos
                \[
                    |f(x)| = |f(x) - L + L| \leq |f(x)-L| + |L| < 1 + |L|.
                \]
            Se $a \notin A$, então $V^*_{\delta}(a) \cap A = V_{\delta}(a) \cap A$ e $f$ é localmente limitada em $a$. Se $a \in A$, então, para todo $x \in V_{\delta}(a) \cap A$, temos
                \[
                    |f(x)| \leq \max \{ 1+|L|, |f(a)| \},
                \]
            de modo que $f$ é localmente limitada em $a$. \blackproof
            \item Seja $\epsilon \in \R_{>0}$. Como $g$ é localmente limitada em $a$, existem $\delta_1, M \in \R_{>0}$ tal que, para todo $x \in A$, se $x \in V_{\delta_1}(a)$, então $|g(x)| \leq M$. Como $\ds \lim_{x \to a} f(x) = 0$, existe $\delta_2 \in \R_{>0}$ tal que, para todo $x \in A$, se $0<|x-a|<\delta_2$, então $|f(x)| < \frac{\epsilon}{M}$. Agora, tome $\delta := \min \{\delta_1, \delta_2\}$. Com isso, para todo $x \in A$, se $0<|x-a|<\delta$, então
                \[
                    |f(x) g(x)| = |f(x)| |g(x)| < M \cdot \dfrac{\epsilon}{M} = \epsilon,
                \]
            de modo que $\ds \lim_{x \to a} f(x) g(x) = 0$. \blackproof
        \end{enumerate}
\end{proof}

\begin{proposition}[Conservação do sinal]
    Sejam $f : A \subseteq \R \to \R$ e $a \in A'$.
        \begin{enumerate}[leftmargin=*, align=left, label=\textbf{(\alph*)}]
            \item Se $\ds \lim_{x \to a} f(x) \in \R_{>0}$, então existe $r \in \R_{>0}$ tal que $f(x) > 0$ para todo $x \in V^*_{r}(a) \cap A$.
            \item Se $\ds \lim_{x \to a} f(x) \in \R_{<0}$, então existe $r \in \R_{>0}$ tal que $f(x) < 0$ para todo $x \in V^*_{r}(a) \cap A$.
        \end{enumerate}
\end{proposition}

\begin{proof}
    Seja $\ds \lim_{x \to a} f(x) = L$.
        \begin{enumerate}[leftmargin=*, align=left, label=\textbf{(\alph*)}]
            \item Como $L \in \R_{>0}$, para $\epsilon = L$ existe $\delta \in \R_{>0}$ tal que, para todo $x \in A$, se $0<|x-a|<\delta$, então $|f(x) - L| < L$; em particular,
                \[
                    0 = L - L < f(x) < 2L.
                \]
            Logo, $f(x) > 0$ para todo $x \in V^*_{\delta}(a) \cap A$. \blackproof
            \item Como $L \in \R_{<0}$, para $\epsilon = -L$ existe $\delta \in \R_{>0}$ tal que, para todo $x \in A$, se $0<|x-a|<\delta$, então $|f(x) - L| < -L$; em particular,
                \[
                    2L < f(x) < 0 = L - L.
                \]
            Logo, $f(x) < 0$ para todo $x \in V^*_{\delta}(a) \cap A$. \blackproof
        \end{enumerate}
\end{proof}

\subsection{Limites Laterais}

\begin{definition}
    Sejam $a \in \R$ e $A \subseteq \R$.
        \begin{enumerate}[leftmargin=*, align=left, label=\textbf{(\alph*)}]
            \item Dizemos que $a$ é um \emph{ponto de acumulação à direita} de $A$ se 
                $
                    \left]a, a+ \delta \right[ \cap A \neq \emptyset
                $
            para todo $\delta \in \R_{>0}$. O conjunto de todos os pontos de acumulação à direita de $A$ é denotado por $A'_+$.
            \item Dizemos que $a$ é um \emph{ponto de acumulação à esquerda} de $A$ se
                $
                    \left]a - \delta, a\right[ \cap A \neq \emptyset
                $
            para todo $\delta \in \R_{>0}$. O conjunto de todos os pontos de acumulação à esquerda de $A$ é denotado por $A'_-$.
        \end{enumerate}
\end{definition}

\begin{definition}[Limites Laterais]
    Seja $f : A \subseteq \R \to \R$.
        \begin{enumerate}[leftmargin=*, align=left, label=\textbf{(\alph*)}]
            \item Dizemos que $L \in \R$ é o \emph{limite lateral à direita} de $f$ em $a \in A'_+$ se para todo $\epsilon \in \R_{>0}$ existe $\delta = \delta(\epsilon, a) \in \R_{>0}$ tal que, para todo $x \in A$, se $a < x < a + \delta$, então $|f(x)-L| < \epsilon$. Isso é denotado por
                \[
                    \lim_{x \to a^+} f(x) = L.
                \]
            \item Dizemos que $L \in \R$ é o \emph{limite lateral à esquerda} de $f$ em $a \in A'_-$ se para todo $\epsilon \in \R_{>0}$ existe $\delta = \delta(\epsilon, a) \in \R_{>0}$ tal que, para todo $x \in A$, se $a - \delta < x < a$, então $|f(x)-L| < \epsilon$. Isso é denotado por
                \[
                    \lim_{x \to a^-} f(x) = L.
                \]
        \end{enumerate}
\end{definition}

\begin{proposition}[Unicidade] \label{prop.calc1:unicidade_limites_laterais}
    Seja $f : A \subseteq \R \to \R$
        \begin{enumerate}[leftmargin=*, align=left, label=\textbf{(\alph*)}]
            \item O limite lateral à direita de $f$ em $a \in A'_+$, quando existe, é único.
            \item O limite lateral à esquerda de $f$ em $a \in A'_-$, quando existe, é único.
        \end{enumerate}
\end{proposition}

\begin{proof}
    \leavevmode
        \begin{enumerate}[leftmargin=*, align=left, label=\textbf{(\alph*)}]
            \item Provemos que se $\ds \lim_{x \to a^+} f(x) = L_1$ e $\ds \lim_{x \to a^+} f(x) = L_2$, então $L_1 = L_2$. Suponha, por absurdo, que $L_1 \neq L_2$, e seja $\epsilon := \dfrac{|L_1 - L_2|}{2} \in \R_{>0}$. Como $\ds \lim_{x \to a^+} f(x) = L_1$, existe $\delta_1 \in \R_{>0}$ tal que, para todo $x \in A$, se $a < x < a+\delta_1$, então $|f(x) - L_1| < \epsilon$. Analogamente, como $\ds \lim_{x \to a^+} f(x) = L_2$, existe $\delta_2 \in \R_{>0}$ tal que, para todo $x \in A$, se $a < x < a+\delta_2$, então $|f(x) - L_2| < \epsilon$. Pois tome $\delta := \min \{\delta_1, \delta_2\}$. Como $a \in A'_+$, existe $x_0 \in A$ tal que $a < x_0 < a+\delta$; com isso, obtemos $|f(x_0) - L_1| < \epsilon$ e $|f(x_0) - L_2| < \epsilon$, donde
                \begin{align*}
                    |L_1 - L_2| &= |(f(x_0) - L_2) - (f(x_0)-L_1)| \\
                    &\leq |f(x_0) - L_2| + |f(x_0)-L_1| \\
                    &< 2\epsilon \\
                    &= |L_1 - L_2|,
                \end{align*}
            uma contradição. Logo, $L_1 = L_2$. \blackproof
            \item Provemos que se $\ds \lim_{x \to a^-} f(x) = L_1$ e $\ds \lim_{x \to a^-} f(x) = L_2$, então $L_1 = L_2$. Suponha, por absurdo, que $L_1 \neq L_2$, e seja $\epsilon := \dfrac{|L_1 - L_2|}{2} \in \R_{>0}$. Como $\ds \lim_{x \to a^-} f(x) = L_1$, existe $\delta_1 \in \R_{>0}$ tal que, para todo $x \in A$, se $a-\delta_1 < x < a$, então $|f(x) - L_1| < \epsilon$. Analogamente, como $\ds \lim_{x \to a^-} f(x) = L_2$, existe $\delta_2 \in \R_{>0}$ tal que, para todo $x \in A$, se $a-\delta_2 < x < a$, então $|f(x) - L_2| < \epsilon$. Pois tome $\delta := \min \{\delta_1, \delta_2\}$. Como $a \in A'_-$, existe $x_0 \in A$ tal que $a-\delta < x_0 < a$; com isso, obtemos $|f(x_0) - L_1| < \epsilon$ e $|f(x_0) - L_2| < \epsilon$, donde
                \begin{align*}
                    |L_1 - L_2| &= |(f(x_0) - L_2) - (f(x_0)-L_1)| \\
                    &\leq |f(x_0) - L_2| + |f(x_0)-L_1| \\
                    &< 2\epsilon \\
                    &= |L_1 - L_2|,
                \end{align*}
            uma contradição. Logo, $L_1 = L_2$. \blackproof
        \end{enumerate}
\end{proof}

\begin{theorem}[Bilateral $\Leftrightarrow$ Laterais] \label{teo.calc1:bilatsselat}
    Sejam $f : A \subseteq \R \to \R$, $a \in A'_+ \cap A'_-$ e $L \in \R$. Então, $\ds \lim_{x \to a} f(x) = L$ se, e somente se, $\ds \lim_{x \to a^+} f(x) = L$ e $\ds \lim_{x \to a^-} f(x) = L$.
\end{theorem}

\begin{proof}
    \leavevmode
        \begin{itemize}[align=left]
            \item[($\Rightarrow$)] Suponha que $\ds \lim_{x \to a} f(x) = L$. Seja $\epsilon \in \R_{>0}$. Como $\ds \lim_{x \to a} f(x) = L$, existe $\delta \in \R_{>0}$ tal que, para todo $x \in A$, se $0<|x-a|<\delta$, então $|f(x)-L| < \epsilon$. Com isso, para todo $x \in A$, se $a < x < a+\delta$, então $0 < x-a < \delta$, isto é, $0 < |x-a| < \delta$, donde $|f(x)-L| < \epsilon$. Logo, $\ds \lim_{x \to a^+} f(x) = L$. Analogamente, para todo $x \in A$, se $a-\delta < x < a$, temos $0 < a-x < \delta$, isto é, $0 < |x-a| < \delta$, donde $|f(x)-L| < \epsilon$. Logo, $\ds \lim_{x \to a^-} f(x) = L$.
            \item[($\Leftarrow$)] Suponha que $\ds \lim_{x \to a^+} f(x) = L$ e $\ds \lim_{x \to a^-} f(x) = L$. Seja $\epsilon \in \R_{>0}$. Como $\ds \lim_{x \to a^+} f(x) = L$, existe $\delta_1 \in \R_{>0}$ tal que, para todo $x \in A$, se $a < x < a+\delta_1$, então $|f(x)-L| < \epsilon$. Como $\ds \lim_{x \to a^-} f(x) = L$, existe $\delta_2 \in \R_{>0}$ tal que, para todo $x \in A$, se $a-\delta_2 < x < a$, então $|f(x)-L| < \epsilon$. Pois tome $\delta := \min\{\delta_1, \delta_2\}$. Com isso, para todo $x \in A$ tal que $0 < |x-a| < \delta$, temos duas possibilidades: ou $x > a$ ou $x < a$. Se $x > a$, então $a < x < a+\delta \leq a+\delta_1$, donde $|f(x)-L| < \epsilon$. Se $x < a$, então $a-\delta_2 \leq a-\delta < x < a$, donde $|f(x)-L| < \epsilon$. Em ambos os casos, $|f(x) - L| < \epsilon$. Logo, $\ds \lim_{x \to a} f(x) = L$.
        \end{itemize}
    Logo, $\ds \lim_{x \to a} f(x) = L \Leftrightarrow \lim_{x \to a^+} f(x) = \lim_{x \to a^-} f(x) = L$. \blackproof
\end{proof}

\begin{proposition}
    Sejam $a \in \R$ e $A \subseteq \R$.
        \begin{enumerate}[leftmargin=*, align=left, label=\textbf{(\alph*)}]
            \item Temos $a \in A'_+$ se, e somente se, $a \in (A \cap \left]a,+\infty \right[)'$.
            \item Temos $a \in A'_-$ se, e somente se, $a \in (A \cap \left]-\infty, a \right[)'$.
        \end{enumerate}
\end{proposition}

\begin{proof}
    \leavevmode
        \begin{enumerate}[leftmargin=*, align=left, label=\textbf{(\alph*)}]
            \item Seja $\delta \in \R_{>0}$.
                \begin{itemize}
                    \item[($\Rightarrow$)] Como $a \in A'_+$, temos $\left]a, a+\delta\right[ \cap A \neq \emptyset$, isto é, existe $x \in \left]a, a+\delta\right[ \cap A$. Como $x \in \left]a, a+\delta\right[$, temos $a < x < a+\delta$, donde $x > a$ e $0 < |x-a| < \delta$. Com isso, $x \in \left]a,+\infty\right[$ e $x \in V^*_{\delta}(a)$. Como $x \in A$, temos $x \in V^*_{\delta}(a) \cap (A \cap \left]a,+\infty\right[)$, de modo que $V^*_{\delta}(a) \cap (A \cap \left]a,+\infty\right[) \neq \emptyset$. Logo, $a \in (A \cap \left]a,+\infty\right[)'$.
                    \item[($\Leftarrow$)] Como $a \in (A \cap \left]a,+\infty\right[)'$, temos $V^*_{\delta}(a) \cap (A \cap \left]a,+\infty\right[) \neq \emptyset$, isto é, existe $x \in V^*_{\delta}(a) \cap (A \cap \left]a,+\infty\right[)$. Como $x \in \left]a,+\infty\right[$, temos $x > a$. Como $x \in V^*_{\delta}(a)$, temos $0 < |x-a| < \delta$, isto é, $a - \delta < x < a + \delta$. Logo, $a < x < a + \delta$, donde $x \in \left]a, a+\delta\right[$. Como $x \in A$, temos $x \in \left]a, a+\delta\right[ \cap A$, de modo que $\left]a, a+\delta\right[ \cap A \neq \emptyset$. Logo, $a \in A'_+$.
                \end{itemize}
            Logo, $a \in A'_+ \Leftrightarrow a \in (A \cap \left]a,+\infty \right[)'$. \blackproof
            \item Seja $\delta \in \R_{>0}$.
                \begin{itemize}
                    \item[($\Rightarrow$)] Como $a \in A'_-$, temos $\left]a-\delta, a\right[ \cap A \neq \emptyset$, isto é, existe $x \in \left]a-\delta, a\right[ \cap A$. Como $x \in \left]a-\delta, a\right[$, temos $a-\delta < x < a$, donde $x < a$ e $0 < |x-a| < \delta$. Com isso, $x \in \left]-\infty,a\right[$ e $x \in V^*_{\delta}(a)$. Como $x \in A$, temos $x \in V^*_{\delta}(a) \cap (A \cap \left]-\infty,a\right[)$, de modo que $V^*_{\delta}(a) \cap (A \cap \left]-\infty,a\right[) \neq \emptyset$. Logo, $a \in (A \cap \left]-\infty,a\right[)'$.
                    \item[($\Leftarrow$)] Como $a \in (A \cap \left]-\infty,a\right[)'$, temos $V^*_{\delta}(a) \cap (A \cap \left]-\infty,a\right[) \neq \emptyset$, isto é, existe $x \in V^*_{\delta}(a) \cap (A \cap \left]-\infty,a\right[)$. Como $x \in \left]-\infty,a\right[$, temos $x < a$. Como $x \in V^*_{\delta}(a)$, temos $0 < |x-a| < \delta$, isto é, $a - \delta < x < a + \delta$. Logo, $a-\delta < x < a$, donde $x \in \left]a-\delta, a\right[$. Como $x \in A$, temos $x \in \left]a-\delta, a\right[ \cap A$, de modo que $\left]a-\delta, a\right[ \cap A \neq \emptyset$. Logo, $a \in A'_-$.
                \end{itemize}
            Logo, $a \in A'_- \Leftrightarrow a \in (A \cap \left]-\infty, a \right[)'$. \blackproof
        \end{enumerate}
\end{proof}

\begin{proposition}
    Seja $f : A \subseteq \R \to \R$ e $L \in \R$.
        \begin{enumerate}[leftmargin=*, align=left, label=\textbf{(\alph*)}]
            \item Se $a \in A'_+$, então
                \[
                    \lim_{x \to a^+} f(x) = L \Leftrightarrow \lim_{x \to a} f \restriction_{A \cap \left]a,+\infty \right[ } (x) = L.
                \]
            \item Se $a \in A'_-$, então
                \[
                    \lim_{x \to a^-} f(x) = L \Leftrightarrow \lim_{x \to a} f \restriction_{A \cap \left]-\infty, a \right[ } (x) = L.
                \]
        \end{enumerate}
\end{proposition}

\begin{proof}
    \leavevmode
        \begin{enumerate}[leftmargin=*, align=left, label=\textbf{(\alph*)}]
            \item Seja $\epsilon \in \R_{>0}$. Como $a \in A'_+$, temos $a \in (A \cap \left]a,+\infty\right[)'$.
                \begin{itemize}
                    \item[($\Rightarrow$)] Como $\ds \lim_{x \to a^+} f(x) = L$, existe $\delta \in \R_{>0}$ tal que, para todo $x \in A$, se $a < x < a+\delta$, então $|f(x) - L| < \epsilon$. Seja $x \in A \cap \left]a,+\infty\right[$ tal que $0 < |x-a| < \delta$. Como $x \in \left]a,+\infty\right[$, temos $x > a$, donde $a < x < a+\delta$. Como $x \in A$, segue que $|f(x) - L| < \epsilon$, isto é, $|f \restriction_{A \cap \left]a,+\infty\right[}(x) - L| < \epsilon$. Logo, $\ds \lim_{x \to a} f \restriction_{A \cap \left]a,+\infty\right[}(x) = L$.
                    \item[($\Leftarrow$)] Como $\ds \lim_{x \to a} f \restriction_{A \cap \left]a,+\infty\right[}(x) = L$, existe $\delta \in \R_{>0}$ tal que, para todo $x \in A \cap \left]a,+\infty\right[$, se $0 < |x-a| < \delta$, então $|f \restriction_{A \cap \left]a,+\infty\right[}(x) - L| < \epsilon$. Seja $x \in A$ tal que $a < x < a+\delta$. Então $x \in \left]a,+\infty\right[$ e $x \in V^*_{\delta}(a) \cap \left]a,+\infty\right[ \cap A$, donde $|f \restriction_{A \cap \left]a,+\infty\right[}(x) - L| < \epsilon$, isto é, $|f(x) - L| < \epsilon$. Logo, $\ds \lim_{x \to a^+} f(x) = L$.
                \end{itemize}
            Logo, $\ds \lim_{x \to a^+} f(x) = L \Leftrightarrow \lim_{x \to a} f \restriction_{A \cap \left]a,+\infty \right[ } (x) = L$. \blackproof
            \item Seja $\epsilon \in \R_{>0}$. Como $a \in A'_-$, temos $a \in (A \cap \left]-\infty,a\right[)'$.
                \begin{itemize}
                    \item[($\Rightarrow$)] Como $\ds \lim_{x \to a^-} f(x) = L$, existe $\delta \in \R_{>0}$ tal que, para todo $x \in A$, se $a-\delta < x < a$, então $|f(x) - L| < \epsilon$. Seja $x \in A \cap \left]-\infty,a\right[$ tal que $0 < |x-a| < \delta$. Como $x \in \left]-\infty,a\right[$, temos $x < a$, donde $a-\delta < x < a$. Como $x \in A$, segue que $|f(x) - L| < \epsilon$, isto é, $|f \restriction_{A \cap \left]-\infty,a\right[}(x) - L| < \epsilon$. Logo, $\ds \lim_{x \to a} f \restriction_{A \cap \left]-\infty,a\right[}(x) = L$.
                    \item[($\Leftarrow$)] Como $\ds \lim_{x \to a} f \restriction_{A \cap \left]-\infty,a\right[}(x) = L$, existe $\delta \in \R_{>0}$ tal que, para todo $x \in A \cap \left]-\infty,a\right[$, se $0 < |x-a| < \delta$, então $|f \restriction_{A \cap \left]-\infty,a\right[}(x) - L| < \epsilon$. Seja $x \in A$ tal que $a-\delta < x < a$. Então $x \in \left]-\infty,a\right[$ e $x \in V^*_{\delta}(a) \cap \left]-\infty,a\right[ \cap A$, donde $|f \restriction_{A \cap \left]-\infty,a\right[}(x) - L| < \epsilon$, isto é, $|f(x) - L| < \epsilon$. Logo, $\ds \lim_{x \to a^-} f(x) = L$.
                \end{itemize}
            Logo, $\ds \lim_{x \to a^-} f(x) = L \Leftrightarrow \lim_{x \to a} f \restriction_{A \cap \left]-\infty, a \right[ } (x) = L$. \blackproof
        \end{enumerate}
\end{proof}
